\documentclass[aps,prb,twocolumn,nofootinbib,superscriptaddress]{revtex4-2}
\usepackage{amsmath,amssymb,amsthm,mathtools,bm}
\usepackage{booktabs}
\usepackage[colorlinks=true,linkcolor=blue,citecolor=blue,urlcolor=blue]{hyperref}
\usepackage{microtype}
\newcommand{\Id}{\hat{\mathbf{1}}}
\newcommand{\id}{\mathbf{1}}
\newcommand{\Tr}{\operatorname{Tr}}
\newcommand{\E}{\widetilde{\mathcal E}}
\newcommand{\D}{\mathcal D}
\newcommand{\Lind}{\mathcal L}
\hypersetup{pdftitle={Channel relaxation and stationary-state spectral gaps},
 pdfauthor={Asadullah Bhuiyan}}
\begin{document}
\title{\texorpdfstring{Channel relaxation and stationary-state spectral gaps\\
\large Internal discussion note}{Channel relaxation and stationary-state spectral gaps: Internal discussion note}}
\author{Asadullah Bhuiyan}
\noaffiliation
\date{October 5, 2026}
\begin{abstract}
We distinguish the relaxation gap of the trajectory-averaged quantum channel
from the occupation gap of its stationary correlation matrix. For the fixed-order,
perfect-reset circuit, the charge-neutral and parity-even channel gaps follow
from a product of single-particle projectors. The stationary occupation spectrum
is a different object: a mode at occupation $1/2$ need not be a slow dynamical
mode. We give an explicit channel-to-Lindbladian construction that preserves
stationary states and controls the relaxation gap, explain what continuity alone
does not establish, and collect the relevant literature with its assumptions.
A matched hard-wall calculation illustrates the distinction without asserting a
thermodynamic purity-gap closing from finite-size data.
\end{abstract}
\maketitle
\setcounter{tocdepth}{1}
\tableofcontents

\section{Which spectrum is gapped?}
Throughout, one time step is one complete measurement cycle. After discarding
the measurement outcomes and tracing out fresh ancillary modes, let $\E$ be
the physical-layer channel. We first fix an invariant operator sector and
assume relaxation to a unique stationary density matrix in that sector.
If $\mu_\nu$ are its nonstationary eigenvalues, define
\begin{align}
 r_{\rm MB}&=\max_{\nu\ {
m nonstationary}}|\mu_\nu|,
 &g_{\rm MB}&=1-r_{\rm MB},\label{eq:gap}\\
 \kappa_{\rm MB}&=-\log r_{\rm MB},
 &g_{\rm MB}&=1-e^{-\kappa_{\rm MB}}.\label{eq:rate}
\end{align}
The channel gap $g_{\rm MB}$ is dimensionless; $\kappa_{\rm MB}$ is an
asymptotic decay rate per cycle. Earlier project plots labeled the corresponding
two-point rate $\Delta_C$; here we write $\kappa_C$ to keep it distinct from
the correlation-matrix gap $\Delta$. For complex multipliers, the relaxation
gap is the deficit of the largest nonstationary \emph{modulus} from unity,
not $\min_\nu|1-\mu_\nu|$. A multiplier near $-1$ can relax slowly despite
being far from $1$ in the complex plane. Additional unit-modulus modes obstruct
mixing; they cannot be discarded to manufacture a positive gap.

The stationary correlation matrix is instead
\begin{equation}
 (\overline G_{\rm ss})_{ij}
 =\Tr(\hat\rho_{\rm ss}\hat\psi_i^\dagger\hat\psi_j),
 \qquad \operatorname{spec}(\overline G_{\rm ss})=\{n_a\}\subset[0,1].
\end{equation}
We use the following explicit conventions:
\begin{equation}
 \Delta=\min_a|n_a-\tfrac12|,
 \qquad \Delta_{\rm pur}=\min_a|1-2n_a|=2\Delta.
 \label{eq:purity}
\end{equation}
Thus $\Delta_{\rm pur}$ is a \emph{two-point purity gap}, not the purity
$\Tr\hat\rho^2$. For a Gaussian state, the single-particle modular
energies are $\epsilon_a=\log[(1-n_a)/n_a]$, so a crossing of $1/2$
is also a zero of the Gaussian modular spectrum. For a non-Gaussian state,
this identifies the modular spectrum of the Gaussian state with the same
two-point function, not that of the actual many-body density matrix.
Covariance closure alone does not establish Gaussianity.

Neither $n_a$ nor $\epsilon_a$ is a channel eigenvalue. In particular,
an occupation-$1/2$ mode is maximally mixed at the two-point level; it need
not be conserved or dark. A pure Gaussian trajectory has a projector-valued
\emph{full-system} $G$, whereas its subsystem correlation matrix may have
occupations near $1/2$. Averaging trajectories and restricting a pure state
are different operations. An occupation crossing alone also does not prove
algebraic spatial correlations or a dissipative critical point.

\section{Exact reset circuit and its two-point map}
We retain the overcomplete Wannier (OW) construction and physical-layer
conventions of Bhuiyan, Pan, and Jian~\cite{BPJ}. Distinct OW modes need not
be orthogonal, and their number stabilizers need not commute.
Let $\{\hat\psi_i,\hat\psi_j^\dagger\}=\delta_{ij}$ and
\begin{equation}
 \hat\chi_j=\sum_i v_{ji}\hat\psi_i,\quad
 v_j^\dagger v_j=1,\quad \hat N_j=\hat\chi_j^\dagger\hat\chi_j.
\end{equation}
Perfect measurement and feedback reset this mode to $s=0$ or $s=1$:
\begin{align}
 \E_{j,0}(\hat\rho)&=\hat\chi_j\hat\rho\hat\chi_j^\dagger
 +(\Id-\hat N_j)\hat\rho(\Id-\hat N_j),\nonumber\\
 \E_{j,1}(\hat\rho)&=\hat\chi_j^\dagger\hat\rho\hat\chi_j
 +\hat N_j\hat\rho\hat N_j.\label{eq:reset}
\end{align}
With $\D[\hat L](\hat\rho)=\hat L\hat\rho\hat L^\dagger
-\tfrac12\{\hat L^\dagger\hat L,\hat\rho\}$, the identities
$\hat N_j^2=\hat N_j$ and $\hat\chi_j\hat\chi_j^\dagger=\Id-\hat N_j$
give, exactly,
\begin{align}
 \E_{j,0}&=\mathrm{Id}+\D[\hat\chi_j]+\D[\hat N_j],\nonumber\\
 \E_{j,1}&=\mathrm{Id}+\D[\hat\chi_j^\dagger]+\D[\hat N_j].
 \label{eq:idplusL}
\end{align}
This is not a first-order approximation to an exponential. The number
dissipator is the measurement dephasing already present in the reset.

For $P_j=v_jv_j^\dagger$ and $Q_j=\id-P_j$, the adjoint dissipators
acting on $\hat\psi_a^\dagger\hat\psi_b$ give
\begin{align}
 \delta_{\rm gain/loss}\overline G
 &=-\tfrac12(P_j\overline G+\overline GP_j)+sP_j,\nonumber\\
 \delta_{\rm dephas.}\overline G
 &=P_j\overline GP_j-\tfrac12(P_j\overline G+\overline GP_j).
 \label{eq:covD}
\end{align}
Here $\delta$ denotes the contribution of the specified dissipator,
not a finite-time integration. Adding them to $\overline G$ yields
\begin{equation}
 \overline G'=Q_j\overline GQ_j+sP_j.\label{eq:elementaryG}
\end{equation}
Equivalently, $\E_{j,s}-\mathrm{Id}$ generates
$\dot{\overline G}=-P_j\overline G-\overline GP_j+P_j\overline GP_j+sP_j$.
No sampling or Wick factorization is needed for this two-point equation.

The ordered product means composition, with step 1 acting first:
\begin{equation}
 \E=\prod_{j=M}^{1}\E_{j,s_j}
 \equiv\E_{M,s_M}\circ\cdots\circ\E_{1,s_1}.
\end{equation}
Define $A=Q_M\cdots Q_1$ and $R_j=Q_M\cdots Q_{j+1}$, with
$R_M=\id$. Then
\begin{equation}
 \overline G'=A\overline GA^\dagger+B,
 \qquad B=\sum_j s_jR_jP_jR_j^\dagger.\label{eq:affine}
\end{equation}
The steady-state equation is affine,
$\overline G_{\rm ss}=A\overline G_{\rm ss}A^\dagger+B$.
Consequently, a stationary many-body state does \emph{not} require an
eigenvalue 1 of $A$. The eigenvalue 1 belongs to $\E$, not necessarily to $A$.

For deviations from stationarity, column vectorization gives
\begin{equation}
 \operatorname{vec}(\delta G')=(A^*\otimes A)\operatorname{vec}(\delta G).
\end{equation}
If $a_p$ are the eigenvalues of $A$, the covariance multipliers are
$a_pa_q^*$. Thus, with $r=\rho(A)=\max_p|a_p|$,
\begin{equation}
 \boxed{g_C=1-r^2,\qquad\kappa_C=-2\log r.}\label{eq:covgap}
\end{equation}
The square comes from the two-sided action on $G$. No eigenvalue of $A$
is removed in computing $r$; an actual unit-modulus mode gives $g_C=0$.
For $Aw=aw$, the matrix $ww^\dagger$ is a covariance eigenmode with
multiplier $|a|^2$. It is not, in general, an eigenprojector of
$\overline G_{\rm ss}$. Nonnormality further distinguishes right state
perturbations from left, observable relaxation modes.

\section{When is this the many-body gap?}
Covariance closure preserves the Heisenberg-space span of $\Id$ and
$\hat\psi_i^\dagger\hat\psi_j$. Its spectrum is therefore included in
the many-body spectrum. In a mixing sector containing these operators,
\begin{equation}
 r_{\rm MB}\ge r^2,\qquad
 g_{\rm MB}\le g_C,\qquad\kappa_{\rm MB}\le\kappa_C.
 \label{eq:bound}
\end{equation}
A finite upper bound alone does not establish a finite many-body gap:
higher correlations could relax more slowly. This is the general
spectral-inclusion statement, closely related to the closed correlation
hierarchy for quadratic Lindbladians~\cite{BZ}.

For the \emph{particular exact resets} in Eq.~\eqref{eq:reset}, a stronger
argument applies. In a basis containing the measured mode, a reset erases
a single measured leg in an even operator; an occupation pair is replaced
by $s$. A normal-ordered $2p$-point neutral correlation therefore maps to
itself and lower orders. The homogeneous block of one cycle is
\begin{equation}
 H_p=(\wedge^p A)\otimes(\wedge^p A)^*,
 \qquad
 \Lambda_{I,J}=\prod_{i\in I}a_i\prod_{j\in J}a_j^*,
 \label{eq:hierarchy}
\end{equation}
up to vectorization order, with $|I|=|J|=p$. The two index sets may
overlap. Antisymmetry produces the minors of $A$; Cauchy--Binet composes
the minors in the actual reset order. The hierarchy is block triangular,
so lower-order source terms do not alter these eigenvalues. This statement
also follows by Schur triangularization when $A$ is defective.

The $p=0$ block is the constant. If $r<1$, every $p\ge1$ multiplier
satisfies $|\Lambda_{I,J}|\le r^{2p}\le r^2$, and $p=1$ attains $r^2$.
These moments span the charge-neutral operator sector. More generally,
an even monomial has at least two legs and obeys the same bound. Hence
\begin{equation}
 \boxed{g_{q=0}=g_{\rm even}=g_C.}\label{eq:equality}
\end{equation}
Charge neutrality means $[\hat N_F,\hat\rho]=0$ and permits mixtures
of number sectors; it is not strong particle-number conservation under
gain and loss. Physical density matrices are parity even. For the
specified Kraus extension to unrestricted operator space, the odd-sector
spectrum instead reaches modulus $r$, giving $g_{\rm full}=1-r$.
The parity-similarity argument and full sector accounting are recorded
in the companion \href{../channel_gap_proof/channel_gap_proof.pdf}{proof note}.

Equation~\eqref{eq:equality} requires perfect resets and the fixed repeated
schedule. It is not a consequence of two-point closure alone, and must not
be exported without proof to imperfect feedback, additional dephasing, or
randomly averaged schedules. A spectral gap also does not fix a
size-uniform mixing time: nonnormal transients, Jordan factors, and
size-dependent prefactors can matter. Nor does it determine the
purification time of individual conditioned trajectories.

\section{No universal implication from the stationary spectrum}
\subsection{Global depolarization: an exactly solvable channel}
An elementary counterexample separates the two questions. For any density
matrix $\hat\sigma$ and $0<p\le1$, the completely positive trace-preserving
map
\begin{equation}
 \mathcal R_{p,\sigma}(\hat X)
 =(1-p)\hat X+p\Tr(\hat X)\hat\sigma
 \label{eq:replacement}
\end{equation}
has unique stationary state $\hat\sigma$. Every traceless operator has
multiplier $1-p$, so its channel gap is $p$, independent of the spectrum
of $\hat\sigma$. Choosing $\hat\sigma$ maximally mixed gives
$\Delta_{\rm pur}=0$ at a strictly positive dynamical gap. Conversely,
$p\to0$ closes the channel gap without changing the stationary state.
For $\hat\sigma=\Id/d$, define $\mathcal T(\hat X)=\Tr(\hat X)\Id/d$.
Since $\mathcal T^2=\mathcal T$, the associated generator has the exact
solution
\begin{align}
 \Lind&=\gamma(\mathcal T-\mathrm{Id}),\nonumber\\
 e^{t\Lind}&=\mathcal T+e^{-\gamma t}(\mathrm{Id}-\mathcal T).
\end{align}
Its eigenvalues are $0$ once and $-\gamma$ on every traceless operator.
The Lindbladian gap is $\gamma$, and its time-$\tau$ channel gap is
$1-e^{-\gamma\tau}$, while the unique endpoint is maximally mixed.
For $N$ fermionic modes, $d=2^N$ and $\overline G_{\rm ss}=\id/2$,
so $\Delta_{\rm pur}=0$ exactly. This global example is not a locality
theorem.

\subsection{On-site depolarization: a size-independent local example}
For a lattice of qubits, let $\mathcal T_i$ replace site $i$ by
$\Id_i/2$ while leaving the reduced state of all other sites unchanged.
The local generators commute, and
\begin{equation}
 \Lind_{\rm loc}=\gamma\sum_i(\mathcal T_i-\mathrm{Id})
 =\frac{\gamma}{4}\sum_{i,\,a=x,y,z}\D[\hat\sigma_i^a].
\end{equation}
A Pauli string with $w$ nonidentity factors is an eigenoperator with
eigenvalue $-\gamma w$. Thus the unique stationary state is
$\Id/2^N$, the many-body Lindbladian gap is $\gamma$, and
$e^{\tau\Lind_{\rm loc}}$ has channel gap $1-e^{-\gamma\tau}$,
all independent of $N$. This is a genuine bounded-range counterexample,
not a global reset in disguise. It is stated for a tensor-product spin
lattice; no locality-preserving spin--fermion identification is assumed.

A directly fermionic one-mode version is balanced loss and gain,
\begin{equation}
 \Lind_{\rm bal}=\eta\bigl(\D[\hat c]+\D[\hat c^\dagger]\bigr),
 \quad \dot n=\eta(1-2n).
\end{equation}
Its stationary state is $\Id/2$ and its full operator spectrum is
$\{0,-\eta,-\eta,-2\eta\}$. The full gap is $\eta$, or $2\eta$
on the physical parity-even single-mode sector. In either convention
it is strictly positive although $n_{\rm ss}=1/2$. This is the simplest
gain/loss realization of the same distinction.

\subsection{If ``Hamiltonian gap'' is meant literally}
A gapped Hamiltonian also admits a maximally mixed \emph{thermal} state:
take $\hat H=\varepsilon\sum_i\hat c_i^\dagger\hat c_i$, with
$\varepsilon>0$. Its many-body excitation gap is $\varepsilon$,
but at inverse temperature $\beta=0$, $\hat\rho_\beta=\Id/2^N$.
There is no contradiction: an excitation gap is not a purity gap of
a finite-temperature state. Unitary evolution under $\hat H$ does
\emph{not} relax every initial state to that endpoint; it preserves the
density-matrix eigenvalues. The first two examples, in contrast,
really do have unique attractive maximally mixed endpoints.

\subsection{Scope of the counterexamples}

Within Gaussian channel theory, the same freedom is visible in
$G'=AGA^\dagger+B$~\cite{DFP}. For example,
$A=\sqrt q\,\id$, $B=(1-q)G_0$, with $0\le q<1$ and
$0\le G_0\le\id$, is an admissible covariance map with fixed point
$G_0$ and $g_C=1-q$. The source and contraction control different
properties. This is an illustration within that channel class, not a
replacement for the reset-sector proof above.

Locality, symmetry, purity, or detailed balance can impose additional
relations. Such relations require their actual hypotheses; they do not
follow from the word ``gap.'' In particular, results linking rapid local
mixing to clustering~\cite{KE} concern spatial correlations, not simply
whether an eigenvalue of $\overline G_{\rm ss}$ equals $1/2$.

\section{A controlled channel--Lindbladian deformation}
There is a precise version of making each complete cycle probabilistic.
Given a channel $\E$, define
\begin{equation}
 \E_p=(1-p)\mathrm{Id}+p\E,\qquad0<p\le1.\label{eq:lazy}
\end{equation}
Its fixed space is exactly that of $\E$, since
$\E_p-\mathrm{Id}=p(\E-\mathrm{Id})$. Its eigenvalues are
$\mu_\nu(p)=1-p+p\mu_\nu$. If the original mixing-sector gap is $g$, then
\begin{equation}
 |\mu_\nu(p)|\le1-p+p|\mu_\nu|\le1-pg,
 \qquad g(p)\ge pg.\label{eq:lazybound}
\end{equation}
Thus a gap already bounded below is preserved for $p>0$. This is a
quantitative statement, not an appeal to continuity.

Set $p=\gamma\,\mathrm dt$, where $\gamma$ is the attempted-cycle rate.
The continuous-time limit is
\begin{equation}
 \Lind_{\rm cycle}=\gamma(\E-\mathrm{Id}).\label{eq:poissonL}
\end{equation}
If $\E(\hat\rho)=\sum_a\hat K_a\hat\rho\hat K_a^\dagger$, completeness
immediately gives
\begin{equation}
 \Lind_{\rm cycle}=\sum_a\D[\sqrt\gamma\,\hat K_a],\qquad
 e^{t\Lind_{\rm cycle}}
 =e^{-\gamma t}\sum_{n=0}^\infty\frac{(\gamma t)^n}{n!}\E^n.
 \label{eq:poisson}
\end{equation}
This is a legitimate Lindblad generator~\cite{Wolf}: complete cycles occur
at Poisson-distributed times. Its stationary states are unchanged, and
its eigenvalues are $\lambda_\nu=\gamma(\mu_\nu-1)$. Therefore
\begin{equation}
 \delta_{\Lind}=\gamma\min_{\nu\ {
m nonstationary}}
 (1-\operatorname{Re}\mu_\nu)\ge\gamma g.
 \label{eq:Lgap}
\end{equation}
For our covariance map, $r^2$ is itself a positive real multiplier and
all other moduli are at most $r^2$. Consequently,
\begin{equation}
 g_C(p)=p(1-r^2),\qquad \delta_{\Lind,C}=\gamma(1-r^2).
\end{equation}
The same equalities hold in the neutral/even many-body sector under
Eq.~\eqref{eq:equality}. The Poissonized rate is not
$\gamma\kappa_C$: random cycle counts and deterministic cycle counts
have different asymptotic rates.

The vanishing of $g(p)$ as $p\to0$ merely reflects a vanishing chance of
updating per discrete step. Measured in physical time with
$p=\gamma\,\mathrm dt$, the generator gap stays finite. At exactly
$p=0$, the discrete channel is the identity and its fixed space enlarges;
that endpoint must not be called a mixing channel.

Conversely, if a relaxing Lindbladian has gap $\delta_{\Lind}>0$, its
finite-time channel $\E_\tau=e^{\tau\Lind}$ satisfies
\begin{equation}
 g(\tau)=1-e^{-\tau\delta_{\Lind}},\qquad\tau>0,
\end{equation}
in the same sector. This does not mean an arbitrary prescribed channel
is $e^{\tau\Lind}$ for some time-independent Lindbladian. Semigroup
embeddability is a separate restriction~\cite{PRZ}; the construction
$\gamma(\E-\mathrm{Id})$ does not assert that $\log\E$ is Lindbladian.

\subsection{Two limitations that matter here}
First, Poissonizing a \emph{whole raster cycle} need not produce a local
sum of bounded-range generators. Its Kraus operators are full-cycle
products, and the sequential depth grows with system size. Instead
weakening each local reset gives
\begin{equation}
 \begin{aligned}
 \prod_{j=M}^{1}\left[\mathrm{Id}
 +\gamma_j\,\mathrm dt(\E_j-\mathrm{Id})\right]\\
 =\mathrm{Id}+\mathrm dt\sum_j\gamma_j(\E_j-\mathrm{Id})
 +O(\mathrm dt^2).
 \end{aligned}
\end{equation}
The resulting local sum generally has a different stationary state from
the ordered finite-reset cycle. A common fixed state of every $\E_j$
would remove this particular obstruction, but must be demonstrated.
The noncommuting OW stabilizers make that distinction important.

Second, continuity of a family of channels or generators alone cannot
exclude gap closure along the path, or in the thermodynamic limit.
To transfer a size-independent gap using Eq.~\eqref{eq:lazybound}, one
needs $\inf_N g_N>0$ and $p$ bounded away from zero, or a fixed physical
rate in the continuous-time limit. Neither a positive finite-size gap
nor a smooth deformation proves this uniform bound. The construction
therefore does not bypass locality-based obstructions to preparing pure
topological states, nor does it establish a gap for a different local
Lindblad approximation to our circuit.

\section{Matched hard-wall calculation}
We compare the exact ordered-product channel gap with the stationary
two-point purity gap for $N_x=20$, $N_y=20,40,60,80,100$, and
$\alpha_1=1,3$. Both calculations use $\alpha_2=30$,
$n_{\rm shell}=1$, inclusive walls $x=5,15$, hard-wall support
truncation, all slabs active, periodic boundaries, zero twist, $X$ trial
orbitals, complex128, and perfect correction including measurement
dephasing. The raster-$y$ schedule uses $x$ outermost, $y$ innermost,
and $A+,A-,B+,B-$ within each cell.

The canonical CPU method
\texttt{classA\_U1FGTN.run\_markov\_channel} evolves the initially
maximally mixed correlation matrix for 60 cycles plus one comparison
cycle. This is exact outcome-averaged two-point evolution, not a finite
sample average. All ten cases pass the recorded stationarity checks:
the last five absolute Frobenius increments are below $10^{-10}$,
and the purity gaps at cycles 60 and 61 differ by less than
$2\times10^{-10}$. The largest final increment is
$6.51\times10^{-15}$. These are numerical convergence checks, not a
rigorous a posteriori bound on fixed-point error for a nonnormal map.

\begin{table}[t]
\caption{Dimensionless channel gap and full-system two-point purity gap
at cycle 61. $N_x=20$ is fixed. The displayed digits summarize
deterministic calculations, not statistical error bars.}
\label{tab:results}
\begin{ruledtabular}
\begin{tabular}{rcccc}
 &\multicolumn{2}{c}{$\alpha_1=1$}&\multicolumn{2}{c}{$\alpha_1=3$}\\
 $N_y$&$g_C$&$\Delta_{\rm pur}$&$g_C$&$\Delta_{\rm pur}$\\
% Generated by build_note.py from verified results.
20 & 0.851086 & 0.086790 & 0.961228 & 0.985595 \\
40 & 0.841480 & 0.004357 & 0.959010 & 0.985659 \\
60 & 0.838258 & 0.021813 & 0.958243 & 0.985667 \\
80 & 0.836644 & 0.005699 & 0.957852 & 0.985668 \\
100 & 0.835675 & 0.009523 & 0.957611 & 0.985668 \\

\end{tabular}
\end{ruledtabular}
\end{table}

Table~\ref{tab:results} shows a large relaxation gap for both alphas,
despite sharply different stationary occupation gaps. At $\alpha_1=1$,
the latter is small and nonmonotonic over the tested sizes; at
$\alpha_1=3$, it remains near $0.986$. These data support a separation
of the two spectral diagnostics. They do \emph{not} yet prove that
$\Delta_{\rm pur}$ vanishes as $N_y\to\infty$, nor a nonzero limiting
channel gap. Furthermore, fixed $N_x=20$ is a fixed-width limit, not a
two-dimensional thermodynamic limit.

The data also contain a separately postprocessed translation average
along $y$. Its momentum-resolved occupations describe that twirled
matrix, not automatically the actual raster-channel fixed point.
Table~\ref{tab:results} uses the untwirled full spectrum. Endpoint
matrices, source hashes, and checksum-verified completion records are
preserved in the matched campaign; the accompanying build record binds
each row to those files. The existing
\href{../../../matched_markov_lindblad_campaign/steady_purity_gap_v1/results/20261005T184136Z/analysis/channel_and_stationary_purity_gaps.pdf}{comparison plot}
is kept separate from this note.

\section{Annotated literature and scope}
The following references address different parts of the argument.
The linked titles and arXiv identifiers in the bibliography provide
direct access; none should be cited as a theorem beyond its hypotheses.

\paragraph{Dissipative topology and the central counterexample.}
Bardyn \emph{et al.}, Ref.~\cite{Bardyn}, Sec.~3.3 and Sec.~7.2,
especially Fig.~6, explicitly separate the purity and dissipative gaps.
Their example has a bulk purity-gap closing while the bulk damping gap
stays finite. This is the most direct precedent for our distinction.
It concerns engineered Gaussian Lindblad dynamics, not the present
ordered reset circuit. The qualifier ``bulk'' matters: open boundaries
can support zero-damping Majorana modes, so the example must not be
described as proving a gapped full open-boundary generator.

\paragraph{Damping, noise, and nonequilibrium criticality.}
Mittal \emph{et al.}, Ref.~\cite{Mittal}, distinguish damping, noise, and
purity scales [Eqs.~(38)--(43)]. In the commuting Gaussian setting,
the stationary covariance involves their ratio. Their special
frustration-free/dark-state conditions lock noise and damping and
give pure stationary states; they do not impose a general identity
between purity and relaxation gaps. This is useful for explaining
which extra structure can restore relations familiar from ground-state
physics, rather than assuming such relations for arbitrary channels.

\paragraph{Symmetry of the generator versus symmetry of its state.}
Mao and Yang, Ref.~\cite{Mao}, relate non-Hermitian symmetry classes
of quadratic Lindblad dynamics to the Altland--Zirnbauer classes of
unique Gaussian stationary states. Their correspondence concerns symmetry
constraints, not equality of spectral gaps. The construction assumes
quadratic Hamiltonians and linear gain/loss jumps. Our number-dephasing
term is quadratic in fermions, and the cycle is ordered and discrete;
two-point closure alone does not place it within that theorem.

\paragraph{Many-body spectra and closed correlation hierarchies.}
Barthel and Zhang, Ref.~\cite{BZ}, Secs.~III.5--III.6, distinguish
unrestricted and two-point relaxation and develop a block-triangular
hierarchy including Hermitian quadratic jumps (Proposition~7).
That is the appropriate continuous-time reference for
Eq.~\eqref{eq:bound}. The exact discrete-reset equality in
Eq.~\eqref{eq:equality} requires our additional higher-moment argument;
it should not be attributed to their covariance equation alone.

\paragraph{Discrete fermionic channels.}
Dierckx, Fannes, and Pogorzelska, Ref.~\cite{DFP}, Sec.~5.1,
give the affine symbol maps and complete-positivity conditions for
fermionic quasi-free channels. This is a direct channel, rather than
only Lindbladian, reference for separating the contraction $A$ from
the source $B$. In the convention of Eq.~\eqref{eq:affine}, the
gauge-invariant Gaussian-channel condition is
$0\le B\le\id-AA^\dagger$. It does not by itself establish
Gaussianity from a closed two-point equation.

\paragraph{Discrete channels with independent spectral diagnostics.}
S\'a \emph{et al.}, Ref.~\cite{Sa}, study random Kraus maps and circuits.
A transition in the support of the complex channel spectrum need not
produce a corresponding stationary-state spectral transition.
This is relevant channel-level evidence for separating the two
diagnostics, but their annulus-to-disk transition is not necessarily
a closing of the relaxation gap, nor a fermionic purity-gap crossing.
Yoshimura and S\'a, Ref.~\cite{YS}, study relaxation in noisy Floquet
circuits, including depolarizing settings with a maximally mixed
stationary state. Their order-of-limits discussion is useful when
comparing weak-noise and thermodynamic limits; it is not a boundary
occupation-gap theorem. Kukulski \emph{et al.}, Ref.~\cite{Kukulski},
provide broader context on random-channel spectra, gaps, and invariant
states, rather than a universal identification of those spectra.

\paragraph{When a dynamical gap constrains spatial correlations.}
Kastoryano and Eisert, Ref.~\cite{KE}, prove clustering statements
under local mixing hypotheses. Their spectral-gap result
(Theorem~9, with a free-fermion corollary) uses regularity and
reversibility assumptions; stronger mixing conditions control further
correlation measures. This is not a theorem that any finite
nonnormal-channel gap forces any chosen stationary occupation gap
to remain open. Those locality and reversibility assumptions have
not been established for the whole-cycle Poissonization here.

\paragraph{Continuous time and embeddability.}
Wolf's notes, Ref.~\cite{Wolf}, provide the standard channel and
Lindblad framework underlying Eq.~\eqref{eq:poisson}. Pucha\l{}a,
Rudnicki, and \.{Z}yczkowski, Ref.~\cite{PRZ}, discuss which Pauli
channels belong to semigroups. They are useful for the warning that
not every channel has a Lindbladian logarithm, not as the source of
our elementary proof that $\gamma(\E-\mathrm{Id})$ is a generator.

\section{Statement suitable for the manuscript}
The relaxation spectrum of the trajectory-averaged channel and the
occupation spectrum of its stationary correlation matrix diagnose
different properties. For the fixed-order perfect-reset circuit, the
parity-even relaxation gap is $g_C=1-\rho(A)^2$, where $A$ is the
ordered product of the complementary OW projectors. The stationary
matrix instead solves an affine fixed-point equation containing the
reset source. An occupation near $1/2$ therefore does not require a
slow dynamical mode. This distinction is consistent with dissipative
topological constructions in which a bulk purity gap closes without
a bulk damping-gap closing~\cite{Bardyn}. Our finite-size hard-wall
data exhibit a small stationary occupation gap at $\alpha_1=1$
alongside a large channel gap; establishing an actual thermodynamic
purity-gap closing requires a separate finite-size analysis.

\begin{thebibliography}{12}
\bibitem{BPJ} A. Bhuiyan, H. Pan, and C.-M. Jian,
\href{https://doi.org/10.1103/nk38-ygyq}{Free-fermion dynamics with measurements:
Topological classification and adaptive preparation of topological states},
Phys. Rev. Research \textbf{8}, 023147 (2026);
\href{https://arxiv.org/abs/2507.13437}{arXiv:2507.13437}.

\bibitem{BZ} T. Barthel and Y. Zhang,
\href{https://doi.org/10.1088/1742-5468/ac8e5c}{Solving quasi-free and quadratic
Lindblad master equations for open fermionic and bosonic systems},
J. Stat. Mech. (2022) 113101;
\href{https://arxiv.org/abs/2112.08344}{arXiv:2112.08344}.

\bibitem{DFP} B. Dierckx, M. Fannes, and M. Pogorzelska,
\href{https://doi.org/10.1063/1.2841326}{Fermionic quasi-free states and maps
in information theory}, J. Math. Phys. \textbf{49}, 032109 (2008);
\href{https://arxiv.org/abs/0709.1061}{arXiv:0709.1061}.

\bibitem{KE} M. J. Kastoryano and J. Eisert,
\href{https://doi.org/10.1063/1.4822481}{Rapid mixing implies exponential
decay of correlations}, J. Math. Phys. \textbf{54}, 102201 (2013);
\href{https://arxiv.org/abs/1303.6304}{arXiv:1303.6304}.

\bibitem{Wolf} M. M. Wolf,
\href{https://mediatum.ub.tum.de/doc/1701036/document.pdf}{Quantum Channels
\& Operations: Guided Tour}, lecture notes (2012), Sec.~7.1.

\bibitem{PRZ} Z. Pucha\l{}a, \L{}. Rudnicki, and K. \.{Z}yczkowski,
\href{https://doi.org/10.1016/j.physleta.2019.04.057}{Pauli semigroups and
unistochastic quantum channels}, Phys. Lett. A (2019);
\href{https://arxiv.org/abs/1903.12448}{arXiv:1903.12448}.

\bibitem{Bardyn} C.-E. Bardyn, M. A. Baranov, C. V. Kraus, E. Rico,
A. Imamoglu, P. Zoller, and S. Diehl,
\href{https://doi.org/10.1088/1367-2630/15/8/085001}{Topology by dissipation},
New J. Phys. \textbf{15}, 085001 (2013);
\href{https://arxiv.org/abs/1302.5135}{arXiv:1302.5135}.

\bibitem{Mittal} R. Mittal, T. Zander, J. Lang, and S. Diehl,
\href{https://doi.org/10.1103/nj2d-l6pz}{Fermion quantum criticality far
from equilibrium}, Phys. Rev. X \textbf{16}, 011069 (2026);
\href{https://arxiv.org/abs/2507.14318}{arXiv:2507.14318}.

\bibitem{Mao} L. Mao and F. Yang,
\href{https://doi.org/10.1103/p8xr-4yzl}{Symmetry classification
correspondence between quadratic Lindbladians and their steady states},
Phys. Rev. B \textbf{112}, 085136 (2025);
\href{https://arxiv.org/abs/2503.10763}{arXiv:2503.10763}.

\bibitem{Sa} L. S\'a, P. Ribeiro, T. Can, and T. Prosen,
\href{https://doi.org/10.1103/PhysRevB.102.134310}{Spectral transitions and
universal steady states in random Kraus maps and circuits},
Phys. Rev. B \textbf{102}, 134310 (2020);
\href{https://arxiv.org/abs/2007.04326}{arXiv:2007.04326}.

\bibitem{YS} T. Yoshimura and L. S\'a,
\href{https://www.nature.com/articles/s41467-024-54164-7}{Robustness of
quantum chaos and anomalous relaxation in open quantum circuits},
Nat. Commun. \textbf{15}, 9808 (2024);
\href{https://arxiv.org/abs/2312.00649}{arXiv:2312.00649}.

\bibitem{Kukulski} R. Kukulski, I. Nechita, \L{}. Pawela, Z. Pucha\l{}a,
and K. \.{Z}yczkowski,
\href{https://doi.org/10.1063/5.0038838}{Generating random quantum channels},
J. Math. Phys. \textbf{62}, 062201 (2021);
\href{https://arxiv.org/abs/2011.02994}{arXiv:2011.02994}.
\end{thebibliography}
\end{document}
